\subsection{The local-field comparison: proved quotient calculations}

\paragraph{Statement D14 (Proved).}\label{res:D14}
Let $K/\mathbb Q_p$ be finite, put
$\Delta=G_{K(\mu_\infty)}$, and write $\pi^{(p')}$ for the maximal
prime-to-$p$ geometric quotient.  Suppose the hyperbolic curve
$U=X\setminus S$ extends to
a smooth proper curve $\mathscr X/\mathcal O_K$ with boundary
$\mathscr S$ finite \'etale over $\mathcal O_K$.
For a point $x\in(\mathscr X\setminus\mathscr S)(\mathcal O_K)$,
\[
       C_{\pi^{(p')}}\bigl(s_x(\Delta)\bigr)=\pi^{(p')}.
\]
For a rational boundary section $c$ choose an integral nonzero tangent
vector $v$, whose reduction is nonzero, and the associated tangential
section $s_v$.  It satisfies the same equality.  For the cuspidal
section $s_a$ obtained from $s_v$ by a Kummer class
$a\in\widehat{K^*}$, let
$b=(b_\ell)_\ell\in\widehat{\mathbb Z}$ be its valuation.  Then
\[
 C_{\pi^{(p')}}\bigl(s_a(\Delta)\bigr)=
 \begin{cases}
       \pi^{(p')},& b_\ell=0\text{ for every }\ell\ne p,\\
       I_c^{(p')},& b_\ell\ne0\text{ for at least one }\ell\ne p.
 \end{cases}                                                   \tag{E.4}
\]
These are fixed groups in the quotient, computed using the induced
actual section action.

\begin{proof}
All finite unramified extensions of $K$ are contained in
$K(\mu_\infty)$: roots of unity of order $q^n-1$, where $q$ is the
residue-field cardinality, give the unramified extension of degree $n$.
Thus $\Delta\subset I_K$.

The specialization input is SGA~1, XIII, Corollary 2.9 and the
construction in 2.10 \cite{D_SGA1}: for a smooth proper morphism with
a relative normal-crossings boundary, the categories of torsors of
order prime to the residue characteristic specialize equivalently.
In the proper case the resulting isomorphism on maximal prime-to-$p$
fundamental groups is also Stacks, Theorem 58.30.3, Tag 0C0R.
Use the section $x$ to identify the fibre functors on the generic and
special fibres.  The resulting based action factors through the
Galois group of the residue field.  Its restriction to $I_K$, hence to
$\Delta$, is the identity automorphism, proving the first equality.
For the integral tangent vector the same comparison uses the tame
local fibre functor: in a parameter with unit tangent, its finite
prime-to-$p$ coverings are obtained by extracting roots of that
parameter of order prime to $p$.  The roots of unity needed for these
coverings are unramified constants.  Specialization therefore
identifies this pointed fibre functor as well.  The local inertia
comparison is explicitly SGA~1, XIII, Lemma 2.11.  It gives the stated
identity action for $s_v$.

Let $\kappa_a:G_K\to\widehat{\mathbb Z}(1)$ be the Kummer cocycle of
the twist.  Since $s_v$ acts trivially on $\pi^{(p')}$ over $\Delta$
and $\chi|_\Delta=1$, the action of $s_a(\sigma)$ for $\sigma\in\Delta$
is the inner automorphism by $\kappa_a^{(p')}(\sigma)\in I_c^{(p')}$.
Every unit of $\mathcal O_K$ has an $n$-th root in $K^{\rm nr}$ when
$n$ is prime to $p$: take a root in the algebraic residue field and
use Hensel's lemma, whose derivative is a unit.  The compatible
division-point sets give the same assertion for the prime-to-$p$
Kummer class of a completed unit.  Consequently on $\Delta$
\[
       \kappa_{a,\ell}=b_\ell\kappa_{\varpi,\ell}\qquad(\ell\ne p),
\]
where $\varpi$ is a uniformizer.
Each $\kappa_{\varpi,\ell}|_\Delta$ is nonzero.  Indeed, the
ramification indices of $K(\mu_\infty)/K$ have bounded
$\ell$-adic valuation for $\ell\ne p$, whereas the value of an
$\ell^n$-th root of $\varpi$ would be $1/\ell^n$ in the
$K$-normalized valuation.  Vanishing of that Kummer class would put
all these roots in $K(\mu_\infty)$, contradicting the bounded
denominators.  Its image is a nonzero closed subgroup of
$\mathbb Z_\ell(1)$, so multiplication by a nonzero $b_\ell$ remains
nonzero.

We record why the geometric centralizer calculation applies to this
quotient.  Finite groups of order prime to $p$ form a class closed
under subgroups, quotients and extensions.  Nakamura
\cite[Lemma (2.1.2)]{Nakamura1994} states for every such full class that
a nonidentity element of a free-generator subgroup has precisely that
subgroup as centralizer.  For at least two punctures a chosen peripheral
loop is part of a free basis.  For one puncture, pass to the compact
curve's mod-$\ell$ homology cover, $\ell\ne p$, as in the last
paragraph of Nakamura's Lemma (2.1.4); it has at least two punctures.
Equivalently, the boundary-separation proof of A2 works over
$\prod_{\ell\ne p}\mathbb Z_\ell$: finite normal covers separate two
cusps and can have at least three boundary points, whose distinct
summand lines in the abelianization intersect trivially.  Nonzero
inertia elements have nonzero positive integral powers in the open
cover group.  Thus distinct inertia groups in the prime-to-$p$
universal cover meet trivially, and the centralizer of every
nonidentity element of $I_c^{(p')}$ is $I_c^{(p')}$.
The image of $\kappa_a^{(p')}$ is either zero or contains such an
element.  This proves both alternatives of (E.4).
\end{proof}

\paragraph{Statement D15 (Open; full and bad-reduction
computations unfinished).}\label{res:D15}
Part E(d) has the good-marked-reduction, prime-to-$p$ calculations D14.
For instance, over $\mathbb Q_5$ on $U_{0,3}$, the section of the
integral point $2$ and the integral tangential section at $0$ have all
of the free prime-to-$5$ group of rank two fixed.  The tangent twist by
$a=5$ has exactly $I_0^{(5')}$ fixed in that quotient.  Hence the
prime-to-$p$ answer already depends on the specified cuspidal section.

The full group is compared with the computed quotient by
$q:\pi\twoheadrightarrow\pi^{(p')}$, with kernel $N$.
There is always the inclusion
\[
       q\bigl(\pi^{s(\Delta)}\bigr)
                    \subseteq(\pi^{(p')})^{s(\Delta)}.          \tag{E.5}
\]
For a fixed quotient element $\bar h$ and a lift $h\in\pi$, the
elements
\[
           d_\sigma=h^{-1}s(\sigma)h s(\sigma)^{-1}\in N
                 \qquad(\sigma\in\Delta)
\]
form a continuous nonabelian $1$-cocycle for the action of $\Delta$ on
$N$.  The element $\bar h$ belongs to the left side of (E.5) exactly
when this cocycle is a coboundary.  Indeed, replacing $h$ by $hn$ makes
it fixed exactly when $d_\sigma=n\,\sigma(n)^{-1}$ for every $\sigma$.
This proves the precise obstruction to passing from D14 to a full-group
calculation.

The smallest retained full-group instance is $U_{0,3}/\mathbb Q_5$
with $s=s_2$, or with the integral tangent at $0$.  The classes just
displayed have not been computed for all quotient elements there.
A retained bad-marked-reduction instance is
$\mathbb P^1\setminus\{0,1,5,\infty\}/\mathbb Q_5$ with the point $2$;
the two boundary sections $0$ and $5$ meet in the special fibre, so
D14's hypothesis fails.  The full and prime-to-$5$ fixed groups for
this instance have not been computed here.

For the requested literature comparison, Mochizuki
\cite[Lemma 1.3.9]{D_Mochizuki2004} assumes two hyperbolic curves over
finite extensions of $p_i$-adic local fields and an isomorphism of
their full arithmetic fundamental groups.  It proves equality of
their types and preservation of cusp inertia subgroups, both in the
full geometric groups and in every pro-$\ell$ quotient.  Its comparison
map is that arithmetic isomorphism.  D14 instead uses the action of a
specified section restricted to $\Delta$; (E.5) gives the available
map between its full and prime-to-$p$ fixed groups.  No full-group
fixed-group formula is inferred from Lemma 1.3.9.  This identifies the
remaining local calculation rather than suppressing it.
